% Pursuing Stacks (cotes 134-2 to 134-8): archive cote 134-4, site page 29 (PDF page 30), typescript page 117.
% typescript. Draft compared: ps3.tex:1428-1492.
% AI-transcribed by Claude Opus 5.5 (claude-opus-5-5, adaptive thinking, effort high), Message Batches API, batch ps2 (msgbatch_01NVwr1Yuo8qJf5rykeGBXPK), 2026-11-12; not yet checked by eye.
\scanpage{134-4/29}
\mspage{117}
\hand{$\Sigma_h$}\xout{xxxf} of intervals in M, satisfying the saturation condition above. The\xout{x}
\xout{xxxxxx} set of all such structures on M will be denoted by $\mathrm{Hom}_3(M)$ , thus
we get two embeddings
\[ \mathrm{Hom}_3(M)\hand{$\hookrightarrow$}\mathrm{Hom}_2(M)\hand{$\hookrightarrow$}\mathrm{Hom}_1(M), \]
in such a way that a weak homotopy interval structure on M may be viewed
also as a particular case of a homotopism structure on M, and a fortiori,
as a particular case of a homotopy relation \unsure{in|on} M. Of course, the homoto-%
py relations or homotopism structures on M we'll ultimately be interested
in, are those stemming from weak homotopy interval structures on M. \hand{Recall}
\over{\illegible{1}}{that} M admits finite products, \inserted{and} these structures satisfy au\over{t}{t}omatically the
accessory assumption, namely commutation of the canonical functor $M\hand{$\to$}M_h$
to finite products.

It is immediate that if we start with a homotopy relation $R_h$ , \over{\illegible{1}}{the}
corresponding \hand{$\Sigma$}$_h$ as defined in A) is saturated. Thus, the canonical
embedding $r_{23}$ of $\mathrm{Hom}_3$ into $\mathrm{Hom}_1$ admits a canonical left inverse $r_{31}$ ,
the restriction of which to $\mathrm{Hom}_2$ is a canonical left inverse $r_{32}$ of the
natural embedding $r_{23}$ of $\mathrm{Hom}_3$ into $\mathrm{Hom}_2$ .

D) \uline{Contractibility structure}. \inserted{(We still assume M admits finite products.)} \hand{(}As a type of structure, it is a
\xout{xxx}set of objects of M
\[ C_h\hand{$\subset$}\mathrm{Ob}(M) , \]
without any ``basic assumption'' on $C_h$ it seems. These objects will be
called the \uline{contractible} objects. Sorry, there \uline{is} a basic assumption
here I \inserted{just} overlooked, \xout{fixxx} namely every X in $C_h$ should have at least one sec-%
tion (thus I better assume beforehand , as in C), that M has a final
object e). To get the natural saturation condition on $C_h$ , \xout{whichxwill}
\over{I'll}{we'll} make use of the third dotted arrow in the organigram, by associa-%
ting to $C_h$ the set \hand{$\Sigma$}$_h$ of ``contractible intervals'', namely intervals
$\underline{I} = (I, \hand{$\delta$}_o,\hand{$\delta$}_1)$ such that I is in $C_h$ . Of course in general there is no
reason that \hand{$\Sigma$}$_h$ should be saturated, never mind - it defines anyhow \inserted{(as seen in C)}
a homotopism notion in M, and hence (as seen in B) a notion of contrac-%
tible objects, i.e. another subset $\overline{C}_h$ of Ob(M). N\over{\illegible{1}}{o}w it occurs to \deleted{\illegible{1}} me
that it is by no means clear that the latter contains $\overline{C}_h$ , which
brings near the necessity of a more stringent \uline{basic assumption} on $C_h$ ,
namely for the very least
\[ C_h\hand{$\subset$}\overline{C}_h \]
(this will imply that any X in $C_h$ has indeed a section over e, as this
is automatically the case for $W_h$-contractible objects). The saturation
condition (Sat 4) will of course be equality
\[ C_h = \overline{C}_h \qquad , \]
\margin{29}
